\documentclass[10pt,letterpaper]{article}

\usepackage[margin=0.68in]{geometry}
\usepackage{array}
\usepackage{booktabs}
\usepackage{longtable}
\usepackage{tabularx}
\usepackage{pdflscape}
\usepackage[table]{xcolor}
\usepackage{titlesec}
\usepackage{fancyhdr}
\usepackage{enumitem}
\usepackage{microtype}
\usepackage[hidelinks]{hyperref}

\definecolor{headingblue}{HTML}{1F4E79}
\definecolor{lightblue}{HTML}{D9EAF7}
\definecolor{softgray}{HTML}{F4F6F8}
\definecolor{darkgray}{HTML}{404040}

\pagestyle{fancy}
\fancyhf{}
\lhead{\small EDCI 59100-016}
\chead{\small Week 4 Lecture Notes}
\rhead{\small Summer 2026}
\cfoot{\small \thepage}
\renewcommand{\headrulewidth}{0.4pt}
\setlength{\headheight}{14pt}

\titleformat{\section}{\normalfont\Large\bfseries\color{headingblue}}{\thesection}{0.55em}{}
\titleformat{\subsection}{\normalfont\large\bfseries\color{headingblue}}{\thesubsection}{0.55em}{}
\titleformat{\subsubsection}{\normalfont\normalsize\bfseries\color{darkgray}}{\thesubsubsection}{0.45em}{}
\titlespacing*{\section}{0pt}{9pt}{4pt}
\titlespacing*{\subsection}{0pt}{7pt}{3pt}
\titlespacing*{\subsubsection}{0pt}{5pt}{2pt}

\setlength{\parindent}{0pt}
\setlength{\parskip}{3.8pt}
\setlist[itemize]{leftmargin=1.2em,itemsep=1.5pt,topsep=2pt}
\setlist[enumerate]{leftmargin=1.35em,itemsep=1.5pt,topsep=2pt}
\renewcommand{\arraystretch}{1.14}
\setlength{\tabcolsep}{4pt}
\emergencystretch=3em

\newcolumntype{L}[1]{>{\raggedright\arraybackslash}p{#1}}
\newcolumntype{Y}{>{\raggedright\arraybackslash}X}

\newcommand{\keyterm}[1]{\textbf{\color{headingblue}#1}}

\begin{document}

\begin{center}
{\Large\bfseries Week 4 Comprehensive Lecture Notes}\\[0.25em]
{\normalsize EDCI 59100-016: Literature Review}\\
{\normalsize Synthesis, Writing Lab Feedback, Thematic Coding, and Draft Review Structure}\\[0.25em]
{\normalsize Md Kamruzzaman Kamrul \textbar{} Summer 2026}
\end{center}

\tableofcontents
\newpage

\section{Week 4 Map and Source Order}

Week 4 continues the earlier literature-review sequence. Week 1 established topic direction, Week 2 built the scoping-review plan, and Week 3 developed article notes and extraction tables. Week 4 turns those separate article notes into synthesis: themes, comparisons, contradictions, gaps, and an organizational plan for the literature review.

The Week 4 materials must be read in this order because some files continue the work from earlier pages:

\begin{enumerate}
\item \textbf{Page 7: Writing Lab Appointment.} This is the Writing Lab reflection requirement and the first major Week 4 task.
\item \textbf{Page 8: Synthesis Matrix and Organizational Outline.} This is the matrix, theme-development, and outline requirement.
\item \textbf{Official Week 4 solution.} This gives the model coding process, draft-literature-review outline, and 50-point Week 4 marking frame.
\item \textbf{Assignment 7 work.} This includes the Purdue Writing Lab email/report, coach feedback, and the submitted reflection.
\item \textbf{Assignment 8 work.} This includes the resource audit, 15-source matrix, themes, and organizational outline.
\item \textbf{Draft and reviewed-draft continuation.} The draft continues from introduction and method into Chapter 3 planning, Chapter 4 architecture themes, and Chapter 5 construction-management applications.
\end{enumerate}

The central intellectual move is from \emph{summary} to \emph{synthesis}. Summary asks what each source says. Synthesis asks what the sources collectively reveal, how they differ, where they contradict, which gaps remain, and how the final paper should be organized.

\section{Writing Lab Appointment and Reflection}

\subsection{Purpose}

The Writing Lab assignment is worth 5 points and is designed to strengthen the literature-review paper through a live meeting with a Purdue Writing Lab coach. The course requires two Writing Lab or coach appointments across the term. The first appointment focuses on topic, patterns, themes, and planned organization. The second appointment focuses on the draft literature-review paper.

The assignment emphasizes that writing coaches help with the writer's needs, not with a fixed script. Appropriate discussion topics include ideation, outline, organization, synthesis across sections, paragraph flow, grammar, spelling, and revision priorities. The instructor also notes that coach fit matters; if one coach is not useful, the student should schedule another coach.

\subsection{Required Process}

To complete the appointment requirement:

\begin{enumerate}
\item Schedule a 45-minute face-to-face or online live appointment with the Purdue Writing Lab.
\item Choose a coach and a time block that gives a full 45-minute session.
\item Complete the appointment form with 3-5 specific issues to discuss.
\item Submit draft paper material or ideas before the appointment.
\item Attend the session and discuss the paper.
\item Submit two items: the coach report/email confirmation and a 1-2 paragraph reflection.
\end{enumerate}

The report does not need to follow a special instructor-facing format. Its function is confirmation and evidence of participation. The reflection should explain what was learned from the session and what revisions will be made.

\subsection{Coach Feedback Actually Received}

The Writing Lab report was from Jacqueline B. after a June 15, 2026 meeting. The session focused on understanding the assignment and methodology. The report noted three main points:

\begin{itemize}
\item The methodology questions were purposeful and well developed.
\item The descriptive analysis should fit a small, emerging body of literature.
\item The discussion/findings should be placed into one chapter rather than repeated across two chapters.
\end{itemize}

The reviewed draft also included two direct notes: the method was drafted according to PRISMA-ScR, and the paper should use two major sections. These comments support a scoping-review framing and a more consolidated findings/discussion structure.

\subsection{Reflection: Method and Structure Decisions}

The main lesson from the Writing Lab session is that Chapter 3 should be a descriptive mapping chapter, not a meta-analysis. Because the topic is new and the studies use varied designs, datasets, models, metrics, and outcomes, the project should chart the field rather than compute pooled effects.

The revision plan should therefore add or strengthen descriptive categories such as number of included studies, publication years, countries or regions, journal types, publication types, study designs, data sources, sample populations, sample sizes, settings, outcomes, model types, statistical or machine-learning methods, and theories or conceptual frameworks. The smaller paper count is defensible because VLM-based construction safety and progress monitoring is an emerging research area. The paper should state this clearly instead of treating it as a weakness.

The second structural lesson is to avoid duplicating findings across two chapters. The better structure is one findings chapter that reports mapped patterns and thematic results in one place, with interpretation integrated after each major finding or moved to a later discussion section without repeating the same evidence.

\section{Synthesis Matrix and Organizational Outline Assignment}

\subsection{Purpose and Required Products}

The Synthesis Matrix and Organizational Outline assignment moves the project from analyzing individual studies to organizing the literature review. It requires three products:

\begin{enumerate}
\item \textbf{Synthesis Matrix, 8 points.} Use approximately 12-15 key sources. Build from the Week 3 Quick Article Notes, but move beyond individual summaries by comparing, grouping, identifying gaps, and adding or revising a theme/category column.
\item \textbf{Theme Development or Categorization Sheet, 6 points.} Identify 2-4 emerging themes. For each theme, provide a label, definition, supporting studies, and notes on tensions, contradictions, or variability.
\item \textbf{Organizational Outline, 6 points.} Create a clear structure for the literature review. Include introduction, review type or methodology, body organization, and a conclusion placeholder.
\end{enumerate}

The assignment allows multiple documents. The matrix and theme sheet can be submitted as the working document in Excel or Google Sheets format, while the outline can be submitted separately in Word or Google Docs format. The naming instruction is: \texttt{Draft Literature Review Plan\_SemYR\_Your LastName\_Your FirstName}.

\subsection{What Counts as Synthesis}

A strong Week 4 matrix demonstrates movement from source-by-source notes into cross-study thinking. The key evidence of synthesis is not volume; it is the ability to show patterns. A strong matrix should:

\begin{itemize}
\item compare similarities and differences among studies;
\item identify contradictions or variability;
\item group studies into emerging themes;
\item mark gaps in the literature;
\item explain why each source matters for the literature-review structure.
\end{itemize}

The assignment explicitly says the matrix does not need to be perfect. What matters is evidence that the writer is thinking across studies and beginning to organize the literature into themes.

\subsection{Rubric Logic}

\begin{longtable}{L{3.0cm} L{9.8cm} L{1.3cm}}
\toprule
\rowcolor{lightblue}
\textbf{Component} & \textbf{Proficient performance} & \textbf{Points}\\
\midrule
\endfirsthead
\toprule
\rowcolor{lightblue}
\textbf{Component} & \textbf{Proficient performance} & \textbf{Points}\\
\midrule
\endhead
\bottomrule
\endfoot
Synthesis Matrix & Clear synthesis across studies; includes themes/categories; shows similarities, differences, contradictions, and appropriate source selection. & 7-8\\
Theme Development & 2-4 clearly defined themes; includes descriptions, supporting studies, and tensions or variability. & 5-6\\
Organizational Outline & Clear, logical, aligned with the chosen review type; supports synthesis rather than listing studies. & 5-6\\
\end{longtable}

Weak submissions read like annotated bibliographies: one row per study with little connection among sources. Strong submissions show how the future review will be structured.

\section{Official Week 4 Model Guidance}

\subsection{Thematic Coding Workflow}

The official Week 4 model presents thematic coding as a four-step process:

\begin{enumerate}
\item \textbf{Familiarization.} Re-read extracted summaries and highlight recurring concepts.
\item \textbf{Open coding.} Assign initial codes such as dataset limitations, temporal attention, field deployment, or evaluation metrics.
\item \textbf{Axial coding.} Group related codes into candidate themes such as modeling approaches, datasets and benchmarks, deployment challenges, and evaluation practices.
\item \textbf{Theme refinement.} Define each theme and list the studies that support it.
\end{enumerate}

This process explains why Week 4 is not only a formatting assignment. It is the stage where extracted evidence becomes the logic of the final literature review.

\subsection{Model Literature-Review Structure}

The official model outline recommends:

\begin{enumerate}
\item Introduction: topic, scope, and research question.
\item Methods: scoping approach, search strategy, and inclusion/exclusion criteria.
\item Descriptive synthesis: corpus characteristics such as year, venue, model type, and dataset.
\item Thematic synthesis: themes supported by evidence and exemplar studies.
\item Discussion: gaps, methodological limits, implications for practice and research.
\item Conclusion: summary and next steps.
\end{enumerate}

The model Week 4 marking frame is broader than Assignment 8's 20-point structure: thematic synthesis quality, 20 points; argument and structure, 10 points; use of evidence, 10 points; writing clarity and formatting, 10 points. The practical implication is that the draft must be evidence-based, structured around the research question, and written as synthesis rather than source listing.

\section{Resource Audit and Corpus Status}

\subsection{Available Evidence}

The local evidence folder contains 73 review-paper PDFs across four groups:

\begin{itemize}
\item baseline papers: sensors, hardware, manual reporting, NLP, and unimodal computer vision;
\item transition papers: transformer action recognition, video mining, synthetic data, and open-vocabulary detection;
\item core papers: VLM, CLIP, image captioning, visual-text similarity, VQA, progress reporting, safety inspection, and activity tracking;
\item gap-analysis and review papers: construction computer vision, safety AI, activity recognition, productivity monitoring, and worker-centric construction scene understanding.
\end{itemize}

Text extraction succeeded for 70 of 73 PDFs. Three PDFs failed full extraction and would require OCR or cleaner copies if exact quotations are needed later:

\begin{itemize}
\item Jung et al. 2021 on 3D CNN real-time action recognition;
\item Zhang et al. 2025 on training-free few-shot tool/material detection;
\item Jeoung et al. 2025 on zero-shot equipment task monitoring.
\end{itemize}

Those three papers do not block Week 4 because two are already represented in the Week 3 source matrix, and Assignment 8 requires synthesis planning rather than final quotation-level full-text verification.

\subsection{Best 15 Sources for the Week 4 Matrix}

\begin{longtable}{L{0.7cm} L{4.0cm} L{4.7cm} L{5.0cm}}
\toprule
\rowcolor{lightblue}
\textbf{No.} & \textbf{Source} & \textbf{Role} & \textbf{Why it matters}\\
\midrule
\endfirsthead
\toprule
\rowcolor{lightblue}
\textbf{No.} & \textbf{Source} & \textbf{Role} & \textbf{Why it matters}\\
\midrule
\endhead
\bottomrule
\endfoot
1 & Zhong et al. (2019) & Construction computer vision review & Establishes the pre-VLM construction CV baseline.\\
2 & Sherafat et al. (2020) & Activity recognition review & Provides worker/equipment temporal-action baseline.\\
3 & Nath et al. (2020) & PPE detection baseline & Shows real-time safety compliance detection.\\
4 & Wu et al. (2021) & CV plus semantic reasoning & Bridges detection and reasoning.\\
5 & Liu et al. (2020) & Image captioning transition & Introduces natural-language scene description.\\
6 & Gil and Lee (2024) & Zero-shot PPE monitoring & Shows PPE monitoring with less annotation dependence.\\
7 & Chen et al. (2024) & AR and vision-language safety query & Shows interactive worker-facing safety support.\\
8 & Wang et al. (2024) & Visual-text safety similarity & Connects site imagery to safety rule language.\\
9 & Chan et al. (2025) & Context-aware ontology-supported VLM agent & Shows domain knowledge and prompt/ontology integration.\\
10 & Hussain et al. (2026) & VLM safety inspection assistant & Represents current assistant-style onsite safety inspection.\\
11 & Bui et al. (2026) & Construction-worker VLM benchmark & Shows strengths and limits of GPT-4o, Florence-2, and LLaVA-1.5.\\
12 & Jung et al. (2024) & VisualSiteDiary & Core progress-reporting and photolog captioning example.\\
13 & Xiao et al. (2024) & Video-to-report generation & Links CV, ChatGPT-style text generation, and daily reports.\\
14 & Jeoung et al. (2025) & Zero-shot equipment task monitoring & Important for spatial, task, and temporal event questions.\\
15 & Zhang et al. (2025) & Training-free few-shot tool/material detection & Important for open-vocabulary and few-shot recognition.\\
\end{longtable}

Optional swaps include Liang et al. (2024) for CLIP/few-shot object recognition, Tohidifar et al. (2024) for synthetic data and annotation, Li et al. (2025) for worker-centric scene understanding, and Rabbi and Jeelani (2024) for broad AI-in-construction-safety review context.

\section{Research Topic, Questions, and Scoping Method}

\subsection{Topic}

The literature-review topic is the use of vision-language models and related multimodal artificial intelligence systems for construction-site safety monitoring and activity tracking. In this context, spatiotemporal reasoning means understanding what is happening, where it is happening, and how that condition changes over time. Examples include PPE use, unsafe worker-equipment proximity, equipment activity, daily progress, and whether a visual scene satisfies a safety rule.

\subsection{Research Questions}

The draft organizes the review around four research questions:

\begin{enumerate}
\item What VLM architectures are currently used for spatiotemporal analysis in construction?
\item How is multimodal data, especially visual and textual data, fused and processed to evaluate safety hazards on active sites?
\item What methodological bottlenecks affect deployment for real-time progress tracking and safety?
\item How are VLMs evaluated and benchmarked for spatiotemporal construction tasks?
\end{enumerate}

\subsection{Review Design}

The review is framed as a PRISMA-ScR scoping review rather than a meta-analysis. This is appropriate because the field is emerging and heterogeneous. The studies vary by task, data source, model type, deployment setting, and evaluation metric. Metrics range from mAP, IoU, precision, recall, and F1 to BLEU, METEOR, CIDEr, SPICE, BERTScore, user ratings, processing time, and accuracy. Because there is no single common effect-size structure, narrative and thematic synthesis are more appropriate than pooled statistical comparison.

\subsection{Search Strategy and Eligibility}

The draft states that the literature search was conducted on March 4, 2026 in Web of Science and Scopus. The search used three conceptual blocks:

\begin{itemize}
\item \textbf{AI architecture:} VLM, vision-text, image-text, VQA, image captioning, CLIP, large multimodal model, foundation model, generative AI, video-language, and video-text.
\item \textbf{Analytical task:} spatiotemporal, video, temporal, dynamic, time-series, 4D, motion, tracking, sequence, action recognition, and activity recognition.
\item \textbf{Application domain:} construction management, civil engineering, AEC, infrastructure, built environment, construction site, safety monitoring, structural health monitoring, and digital twin.
\end{itemize}

Backward snowballing was added because new VLM terminology may miss foundational or transitional work that does not use current labels.

The inclusion criteria emphasize 2019-2026 core VLM literature, 2010-2026 foundational CV baselines, peer-reviewed English journal articles, multimodal architectures integrating visual/geometric and textual/semantic data, and active-construction-phase applications in safety management or progress tracking. Exclusions include pre-2010 items, post-March 2026 items, editorials, book chapters, theses, non-English documents, conference proceedings, industry white papers, purely unimodal methods unless used as baselines, and post-construction/facility-management domains.

\subsection{Selection and Charting}

The draft selection process reports 1,443 raw records: 769 from Web of Science and 674 from Scopus. Automated filters removed 1,320 records, leaving 123 articles. Removing 33 duplicates left 90 unique records. Title/abstract screening excluded 32 records, leaving 58 for full-text review. Full-text screening removed 33 more, yielding 25 core VLM and transition articles. Backward snowballing added 31 baseline and transition articles, producing a final corpus of 56 articles.

The data charting framework has three dimensions:

\begin{longtable}{L{3.0cm} L{4.7cm} L{6.0cm}}
\toprule
\rowcolor{lightblue}
\textbf{Dimension} & \textbf{Variables} & \textbf{Purpose}\\
\midrule
\endfirsthead
\toprule
\rowcolor{lightblue}
\textbf{Dimension} & \textbf{Variables} & \textbf{Purpose}\\
\midrule
\endhead
\bottomrule
\endfoot
Content analysis & Cluster/application, input data, spatiotemporal task, automation level & Identifies what the AI system is doing in construction context.\\
Technical analysis & Algorithm/model, dataset, validation metric, compute/hardware & Explains how the system works and how it is evaluated.\\
Critical analysis & Key limitation, research gap, future direction & Interprets significance, unresolved problems, and next research needs.\\
\end{longtable}

\section{Chapter 3 Planning: Descriptive Mapping}

Chapter 3 should not be a meta-analysis. It should map the included literature. The Writing Lab feedback and draft notes identify the following descriptive-analysis categories:

\begin{itemize}
\item number of studies;
\item publication years and year-by-year trends;
\item countries, regions, journals, publication types, CiteScore, and SJR impact indicators;
\item document co-citation clusters and highly cited articles;
\item top cited authors and collaboration networks by author and institution;
\item keyword frequency, TF-IDF keywords by year, word clouds, keyword co-occurrence, and technology-frequency heatmaps;
\item study characteristics: research design, data sources, sample population, sample size, and setting;
\item topic/content characteristics: outcomes examined and what each study considered;
\item methodological characteristics: modeling approaches, statistical techniques, machine-learning methods, survey methods, logistic models, and cross-sectional designs;
\item conceptual or theoretical grounding, including whether studies lack explicit theory.
\end{itemize}

The key revision principle is to connect every descriptive table or graph back to the research questions. A figure is useful only if it helps explain the development of VLM architectures, multimodal fusion, deployment bottlenecks, or benchmark practices.

\section{Theme Development}

\subsection{Theme 1: From Detection to Semantic Multimodal Reasoning}

\textbf{Definition.} Earlier studies detect visible objects, workers, equipment, PPE, or actions. Newer studies interpret meaning by connecting images and videos to captions, questions, prompts, safety rules, or domain knowledge.

\textbf{Supporting studies.} Zhong et al. (2019), Nath et al. (2020), Wu et al. (2021), Liu et al. (2020), Gil and Lee (2024), Wang et al. (2024), Chan et al. (2025), Hussain et al. (2026), Bui et al. (2026), and Zhang et al. (2025).

\textbf{Synthesis.} The field does not jump directly from computer vision to VLMs. It moves through layers: object detection, semantic reasoning, image captioning, visual-text matching, VQA, ontology-guided agents, and multimodal assistants. Traditional detection is easier to evaluate because object labels and bounding boxes are clear. Semantic multimodal reasoning is more useful for construction decisions, but harder to validate because it is prompt-sensitive, context-dependent, and vulnerable to hallucination or missing domain knowledge.

\subsection{Theme 2: Safety Compliance and Contextual Hazard Interpretation}

\textbf{Definition.} Safety monitoring begins with visible compliance checks such as hardhats and vests, then expands toward worker behavior, unsafe proximity, safety-rule matching, ontology-supported compliance, and assistant-style inspection.

\textbf{Supporting studies.} Nath et al. (2020), Gil and Lee (2024), Chen et al. (2024), Wang et al. (2024), Chan et al. (2025), Hussain et al. (2026), and Bui et al. (2026).

\textbf{Synthesis.} PPE detection is a practical starting point because it is visible and measurable. Gil and Lee reduce annotation dependence through zero-shot monitoring. Wang et al. link visual site evidence to textual safety requirements. Chen et al. and Hussain et al. shift toward interactive assistance. Chan et al. adds ontology and construction-specific knowledge to reduce generic VLM weaknesses. The unresolved issue is that a worker can wear PPE and still be unsafe because risk depends on relationships among worker action, equipment proximity, work sequence, site condition, and safety-rule context.

\subsection{Theme 3: Activity Tracking, Progress Monitoring, and Temporal Reasoning}

\textbf{Definition.} Progress monitoring and activity tracking require models to understand changes over time: worker actions, equipment tasks, daily reports, productivity, and sequence-level site activity.

\textbf{Supporting studies.} Sherafat et al. (2020), Liu et al. (2020), Jung et al. (2024), Xiao et al. (2024), Jeoung et al. (2025), and Bui et al. (2026).

\textbf{Synthesis.} Activity recognition existed before VLMs, but current work links visual evidence to language and reporting. Liu et al. shows early image captioning for construction activities. Jung et al. and Xiao et al. connect images/videos to daily-report generation. Jeoung et al. separates spatial, task, and temporal event questions for equipment monitoring. Bui et al. is important because it shows static images remain weak for dynamic worker behavior. The key gap is the difference between describing frames and reasoning over long-horizon site activity.

\subsection{Theme 4: Dataset, Benchmark, and Deployment Limits}

\textbf{Definition.} Across safety and progress studies, the same limitations recur: small construction-specific datasets, inconsistent metrics, domain shift, occlusion, lighting, clutter, scarce real-site validation, weak shared benchmarks, and sensitivity to prompt design.

\textbf{Supporting studies.} Liu et al. (2020), Davila Delgado and Oyedele (2021), Tohidifar et al. (2024), Xin et al. (2024), Bui et al. (2026), Li et al. (2025), Rabbi and Jeelani (2024), Zhong et al. (2019), Chan et al. (2025), and Xiao et al. (2024).

\textbf{Synthesis.} Data is not a side issue. It shapes what models can do. Zero-shot, few-shot, synthetic-data, and pretrained VLM approaches reduce annotation cost, but they do not remove the need for field validation. Models that work on curated images may fail under occlusion, poor lighting, camera-angle variation, site clutter, and site-specific terminology.

\section{Organizational Outline for the Literature Review}

\subsection{Introduction}

The introduction should define the practical problem: construction sites are dynamic, labor-intensive, equipment-rich, and difficult to monitor through manual observation alone. Manual safety inspection and daily reporting are slow, subjective, and labor-intensive. The review topic is VLMs and related multimodal AI systems for construction safety monitoring and activity tracking. The introduction should define spatiotemporal reasoning plainly and state that the review maps architectures, data modalities, tasks, metrics, deployment barriers, and gaps.

\subsection{Scoping Review Approach}

The methods section should describe the database search, snowballing, inclusion/exclusion criteria, screening process, and charting categories. It should explicitly justify the scoping-review method because the literature is broad and heterogeneous. The section should explain that the review maps patterns rather than calculating effect sizes.

\subsection{Body Organization}

The body should use a thematic and developmental structure:

\begin{enumerate}
\item From detection to semantic multimodal reasoning.
\item Safety compliance and contextual hazard interpretation.
\item Activity tracking, progress monitoring, and temporal reasoning.
\item Datasets, metrics, and field deployment barriers.
\end{enumerate}

This structure supports synthesis because each section compares groups of studies rather than listing one paper at a time.

\subsection{Conclusion Placeholder}

The conclusion should state the main pattern: construction AI research is moving from detecting visible site elements toward interpreting safety and progress meaning through language, context, and multimodal reasoning. The main future directions are construction-specific VLM benchmarks, stronger temporal datasets, transparent field validation, ontology/prompt reliability, and systems that connect visual observations to safety rules and work progress in ways practitioners can trust.

\section{Chapter 4 Continuation: VLM Architecture Framework}

\subsection{Three-Layered Spatiotemporal VLM Framework}

The draft's Chapter 4 shifts from descriptive mapping to thematic analysis of actual architectures. It proposes a three-layer framework based on the automated-systems analogy of sensors, processors, and actuators:

\begin{enumerate}
\item \textbf{Layer 1: Multimodal inputs.} Raw visual streams, text rules, domain knowledge, BIM/IFC data, and kinematic or sensor data.
\item \textbf{Layer 2: Spatiotemporal VLM core.} The processing layer where perception, fusion, reasoning, prompting, captioning, VQA, and decision logic occur.
\item \textbf{Layer 3: Construction-management outputs.} Safety management, risk interpretation, progress tracking, daily reporting, inspection decisions, and automated interventions.
\end{enumerate}

The reason for emphasizing spatiotemporal analysis is that construction sites are dynamic. Safety and progress depend on location, sequence, duration, proximity, and changing site state. Static image recognition can be useful, but construction monitoring ultimately requires reasoning across physical space and time.

\subsection{Paradigm A: Zero/Few-Shot Prompting}

Paradigm A responds to the data-burden problem in traditional supervised computer vision. Construction datasets are expensive to annotate, computationally demanding, and often too narrow for diverse site conditions. Zero-shot and few-shot methods use pretrained models such as CLIP, DINOv2, and Florence-2 to recognize objects or classes through natural-language prompts rather than full task-specific retraining.

The mechanism is a shared embedding space. Visual features and text prompts are compared, often through cosine similarity. This allows models to recognize temporary objects, equipment tasks, tools/materials, and PPE categories by changing the text query. The weakness is semantic expressiveness: these models often output object tags, bounding boxes, or isolated labels, but not full situational understanding. They also lack construction-specific knowledge unless adapted.

\subsection{Paradigm B: Domain Fine-Tuned Captioning and Generative Description}

The draft labels one heading as Paradigm C, but the content describes the second paradigm: generative adaptation through domain fine-tuned image captioning. Image captioning combines a visual encoder with a language decoder to translate site images into natural-language descriptions.

This paradigm matters because it turns raw visual evidence into human-readable narratives. It supports safety narratives, hazard descriptions, photolog captioning, and daily construction reporting. VisualSiteDiary is a key example: it captions site photologs and supports image retrieval for daily reporting.

The weakness is interactivity. A captioning model produces a fixed description and stops. A manager or worker cannot ask follow-up questions, query an unmentioned hazard, or dynamically compare the scene against multi-step safety rules.

\subsection{Paradigm C: Hybrid Geometric-Semantic and Agentic Systems}

Paradigm C combines geometric precision with language-based interpretation. Examples include YOLO or object detection arrays feeding an LLM/MLLM, construction safety ontologies shaping prompts, VQA systems answering compliance questions, AR interfaces supporting worker queries, and robots comparing captured images with BIM/IFC metadata.

The advantage is two-way reasoning. The system can inspect visual evidence, refer to domain rules, produce captions, answer questions, and support safety or progress decisions. The weakness is reliability. These systems remain sensitive to prompt wording, ontology maintenance, environmental noise, device limitations, latency, and complex dynamic behavior.

\section{Chapter 5 Continuation: Construction Applications}

\subsection{Safety and Risk Management}

Construction safety applications begin with PPE compliance because hardhats, vests, and visible worker equipment are easy to define and evaluate. Zero-shot retrieval approaches can compare image captions or visual features with OSHA-related text prompts. However, safety risk is not reducible to PPE. A fully equipped worker may still be at risk because of equipment proximity, fall exposure, unsafe sequence, or surrounding site context.

The literature therefore moves into richer safety interpretation:

\begin{itemize}
\item image captioning plus object detection to describe unsafe behaviors;
\item visual-text similarity to match images with safety-rule language;
\item AR-based Visual Construction Safety Query systems for worker-facing hazard questions;
\item ontology-guided VLM agents to reduce hallucination and enforce domain rules;
\item mobile or assistant-style systems that combine geometric detections with generated hazard captions.
\end{itemize}

The unresolved gap is nuanced behavioral interpretation. General VLMs may struggle to distinguish collaboration from communication, identify subtle unsafe behavior, or reason reliably under cluttered and changing site conditions.

\subsection{Progress Tracking and Automated Reporting}

Progress tracking uses the same spatiotemporal engines but applies them to productivity and documentation. The first need is object and resource localization: tools, materials, equipment, and temporary site objects. Zero-shot and few-shot recognition reduce retraining burdens for these categories.

The next step is report generation. VisualSiteDiary and video-to-report systems convert photologs, images, or video into text for daily construction reports. These systems help automate tedious administrative documentation, but reporting quality depends on dataset coverage, historical context, and language-generation reliability.

The strongest progress-monitoring challenge is time. A static image may show a tool or machine, but it does not show the full sequence of work. Jeoung et al. (2025) is important because it asks spatial, task, and temporal event questions over video frames. Robotic and BIM-linked systems add another layer: a quadruped robot can inspect a partially built environment and compare images against IFC metadata to decide whether architectural elements are present and correctly installed.

\section{Technical Synthesis of Core VLM Applications}

\begin{landscape}
\scriptsize
\rowcolors{2}{softgray}{white}
\begin{longtable}{L{2.9cm} L{2.4cm} L{1.6cm} L{1.8cm} L{2.4cm} L{2.7cm} L{2.2cm} L{4.9cm}}
\caption{Master Data Charting and Technical Synthesis of Core VLM Applications}\\
\toprule
\rowcolor{lightblue}
\textbf{Source} & \textbf{Domain} & \textbf{Paradigm} & \textbf{Input} & \textbf{Task} & \textbf{Model} & \textbf{Metric} & \textbf{Key limitation}\\
\midrule
\endfirsthead
\toprule
\rowcolor{lightblue}
\textbf{Source} & \textbf{Domain} & \textbf{Paradigm} & \textbf{Input} & \textbf{Task} & \textbf{Model} & \textbf{Metric} & \textbf{Key limitation}\\
\midrule
\endhead
\bottomrule
\endfoot
Zhang et al. (2025) & Progress: tools/materials & A: zero-shot & Images & Object recognition & DINOv2, CLIP, TIP-X & mAP & False positives from limited domain-specific architectural knowledge in pretrained models.\\
Gil and Lee (2024) & Safety: PPE & A: zero-shot & Images & Zero-shot monitoring & Image captioning plus cosine similarity & Accuracy, F1, F2 & Struggles with complex behaviors and domain-specific terminology.\\
Liang et al. (2024) & Progress: temporary objects & A: zero-shot & Images & Few-shot detection & CLIP plus multimodal prototypes & Accuracy & Relies on similarity distributions rather than precise instance geometry.\\
Liu et al. (2020) & Safety and progress & B: fine-tuned & Images & Image captioning & CNN-LSTM & BLEU, ROUGE, CIDEr & Generates static descriptions and lacks interactive feedback.\\
Xiao et al. (2022) & Progress: equipment & B: fine-tuned & Images & Object/action recognition & Transformer-SCST & Precision, recall, F1 & Captioning-based object recognition can trail bounding-box detectors.\\
Zhai et al. (2023) & Safety: behaviors & B: fine-tuned & Images & Unsafe behavior captioning & Transformer plus attention & BLEU, precision, recall & Sentence variation is limited and annotation quality matters heavily.\\
Yang et al. (2023) & Safety: action ID & B: fine-tuned & Video & Behavior recognition & STR-Transformer & Loss, accuracy & High computational complexity limits real-time feasibility.\\
Wang et al. (2024) & Safety: hazard ID & B: fine-tuned & Image plus text & Semantic information extraction & Transformer encoder plus USE & Similarity matrix & Complex scenes with sparse descriptions remain difficult.\\
Jung et al. (2024) & Progress: DCRs & B: fine-tuned & Photologs & Captioning and retrieval & VisualSiteDiary, mPLUG & BLEU, CIDEr, SPICE & Flattened grid visual features can lose spatial information.\\
Chen et al. (2024) & Safety: VQA & C: hybrid & Image plus text & VQA/image-text retrieval & VLM plus AR & Recall@5, SPICE, accuracy & Complex danger-type and recommendation queries suffer from linguistic ambiguity.\\
Hussain et al. (2026) & Safety: inspection & C: hybrid & Images & Hazard detection and captioning & YOLO plus LLM & BLEU, ROUGE, BERTScore & Mobile processing limits and dynamic site conditions affect stability.\\
Chan et al. (2025) & Safety: compliance & C: hybrid & Image plus text & VQA/compliance checking & CogAgent plus LoRA plus ontology & Sensitivity, specificity, F1 & Ontology must be manually updated for site-specific regulations.\\
Xiao et al. (2024) & Progress: daily logs & C: hybrid & Video & Activity recognition plus report generation & YOLO, 3D ResNet, ChatGPT & Likert/user experience & Sparse historical data can produce hallucinated or inconsistent report text.\\
Jeoung et al. (2025) & Progress: task monitoring & C: hybrid & Video frames & Temporal event analysis & Florence-2, SAM2, GPT-4o & Accuracy, time & Complex temporal questions show lower accuracy and prompt sensitivity.\\
Jang et al. (2025) & Progress: action recognition & C: hybrid & Video & Action recognition & TimeSformer & Accuracy, F1 & Visually similar installation actions can be confused.\\
Tuomisto et al. (2026) & Progress: Scan-to-BIM & C: hybrid & Images plus BIM & Object presence detection & Claude 3.7 Sonnet plus robot & Accuracy, precision & Requires high LOD 300 BIM accuracy and struggles with multi-stage elements.\\
\end{longtable}
\end{landscape}

\section{Final Week 4 Takeaways}

\begin{enumerate}
\item Week 4 is the point where the project becomes a literature review rather than a set of article notes.
\item The Writing Lab feedback clarifies that the method should be descriptive, mapping-oriented, and PRISMA-ScR aligned.
\item The smaller number of papers is defensible because the VLM construction-monitoring field is emerging.
\item Assignment 8 requires three connected products: synthesis matrix, theme-development sheet, and organizational outline.
\item The strongest structure is thematic: detection to multimodal reasoning, safety hazard interpretation, progress and temporal reasoning, and dataset, benchmark, and deployment limits.
\item Chapter 3 should map the corpus; Chapter 4 should explain VLM architectures; Chapter 5 should apply those architectures to safety and progress.
\item The main technical evolution is from object detection to captioning to interactive, ontology-guided, agentic VLM systems.
\item The main gap is still reliability in real construction environments: long-horizon temporal reasoning, clutter, occlusion, prompt sensitivity, domain knowledge, benchmarks, and field validation.
\end{enumerate}

\section*{Core References Mentioned in Week 4}

\small
Alaloul et al. (2022); Asadzadeh et al. (2020); Bohn and Teizer (2010); Bui, Chodosh, and Tavakoli (2026); Chan et al. (2025); Chen et al. (2023, 2024); Davila Delgado and Oyedele (2021); Gil and Lee (2024); Hussain et al. (2026); Jeoung et al. (2025); Jiang and Messner (2023); Jung et al. (2024); Li et al. (2025); Liang et al. (2024); Liu et al. (2020); Nath, Behzadan, and Paal (2020); Paneru and Jeelani (2021); Rabbi and Jeelani (2024); Samsami (2024); Sherafat et al. (2020); Tohidifar et al. (2024); Tricco et al. (2018); Tuomisto et al. (2026); Wang and El-Gohary (2023, 2024); Wang et al. (2024); Wu et al. (2021); Xiao et al. (2022, 2024); Yang et al. (2023); Zhang and El-Gohary (2015); Zhang et al. (2025); Zhong et al. (2019).

\end{document}
